\begin{center}
\begin{tabular}{@{}lrr@{}}
\toprule
Fresh literal check & AX11QGH & AX11HR\\\midrule
Working states / transition rows &194 / 388&194 / 388\\
Primitive source sites &3,671&3,671\\
Symbolic outcomes &4,778&4,778\\
Generic list lemmas derived from literal rows &4&4\\
Independent interpreter events &1,106,693&1,013,419\\
Maximum events for one outcome &385&373\\
Minimum constant in exact physical duration &3&3\\
Complete blank startup transitions &57,922&48,320\\
Independent source resolvers agree &yes&yes\\
Arithmetic-box oracle used by literal replay &no&no\\
Dense/native table conjugacy &yes&yes\\
Distinct from all 110 cataloged 194-state classes &yes&yes\\
End-to-end Lean/Rocq kernel verification &not supplied&not supplied\\
\bottomrule
\end{tabular}
\end{center}
An interpreter event can stand for an arbitrary-length unary or list sweep;
its count is not an upper bound on physical Turing steps. The generic sweep
identities and positive affine physical durations are what justify the
universal arithmetic conclusion. The scopes include nine independent
arbitrary natural register values, both decrement branches, arbitrary
extra gap, arbitrary nonempty finite spare list, and complete blank
exteriors.

Novelty was compared by exact canonical transition tuples against 175
spanning presentations covering all 110 catalog classes. Two nonspanning
diagnostic tables are also recorded and excluded by the spanning-root
invariant. Neither candidate matches any input under arbitrary control
renaming, any spanning root, global tape reflection, global symbol
complement, or omission of the write/move on terminal edges. Symbol
complement here is an extra structural test; it is not claimed to preserve
zero-blank execution. The two new candidates also have different such
canonical tuples from each other.
